% ===================================================================
%  TABELO N° 16 — La Granda Magazino. — La Bazaro. La Ludili.
%  Feuillets 103 a 108 du fac-simile = folios 101 a 106.
%  Correspondance : folio imprime = numero de feuillet - 2.
%
%  Ca esas la dernier tabelo dil livreto. Le feuillet 109 (« TABELO
%  DIL KONTENAJO ») ne aparas ci : olu ne esas parto dil tabelo 16.
%
%  Le sub-titulo « La Granda Magazino. --- La Bazaro. » e « La
%  Ludili. » esas komposita en petit-kapitali (\VUpk), NE en grase
%  quale por la tabeli 1-9 : lo esas fidela a l'aspekto dil fac-
%  simile, ne un eroro di transskripto.
% ===================================================================

% --- Feuillet 103 (folio 101 : ouverture de tableau) ----------------
\begin{VUpage}[103]{101}
%%K t16-apar-1 apar
\VUsaut{4.0mm}
\VUtitre{15.0pt}{60}{TABELO N\textsuperscript{o} 16}
\VUsaut{2.0mm}
\VUpk{13.2pt}{40}{La Granda Magazino. --- La Bazaro.}
\VUpk{13.2pt}{40}{La Ludili.}
\VUsaut{3.4mm}

%%K t16-01-1 p
\VUblancAlinea
1. --- La nov yaro proximeskas. La granda\nl
magazini preparis sua estaleyi de \VUgras{ludili} \textsuperscript{(1)} e\nl
novyaral donaci.

%%K t16-01-2 p
\VUblancAlinea
Me demandis de mea avulo, ke il duktez me\nl
aden la bazaro por vidar ta bel expozeyo, ed il\nl
konsentis lo.

%%K t16-02-1 p
\VUblancAlinea
2. --- Unesme ni haltis avan bel armo-skudi,\nl
qui kontenas \VUgras{fusilo}, \VUgras{kepio}, \VUgras{sabro}, \VUgras{espado-zono}\nl
e paro de \VUgras{epoleti}, o \VUgras{kasko}, \VUgras{kuraso} \textsuperscript{(3)} e \VUgras{pisto}\cc
\VUgras{lo} \textsuperscript{(4)}. To omna interesis me multe plu kam la\nl
\VUgras{lum-projektili} \textsuperscript{(2)}.

%%K t16-tit-1 sub
\VUsaut{3.0mm}
\VUpk{10.2pt}{40}{Bela vidaji.}
\VUsaut{2.6mm}

%%K t16-03-1 p
\VUblancAlinea
3. --- En la sama fako esis \VUgras{facionero} \textsuperscript{(5)} en\nl
sua \VUgras{garito} ; il semblis facionar avan tota mena\cc
jerio kartona. Quante multa animali esis ibe :\nl
\VUgras{Papagayo} \textsuperscript{(6)} sur sua perchilo, \VUgras{rinocero} \textsuperscript{(7)} kun\nl
korno sur la nazo, \VUgras{rentiro} \textsuperscript{(8)}, \VUgras{strucho} \textsuperscript{(9)}, \VUgras{si}\cc
\VUgras{mio} \textsuperscript{(10)} krokanta nuco, \VUgras{foxo} \textsuperscript{(11)} forportanta\nl
hanino, \VUgras{krokodilo} \textsuperscript{(12)} beanta per enorma\nl
maxili, \VUgras{kamelo} \textsuperscript{(13)}, eleganta \VUgras{levriero} \textsuperscript{(14)}, \VUgras{pan}\cc
\VUgras{tero} \textsuperscript{(15)}, pose du bestii quin me ne konocis e\nl
qui esas, segun dico, \VUgras{vizelo} \textsuperscript{(16)} e \VUgras{hieno} \textsuperscript{(17)}.\nl
Anke ibe esis \VUgras{leopardo} \textsuperscript{(18)} e \VUgras{volfo} \textsuperscript{(21)}; tote\nl
proxime esis miniona \VUgras{rateti} \textsuperscript{(19)} e mikrega \VUgras{mu}\cc
\VUgras{seti} \textsuperscript{(20)} kauchuka. Laste \VUgras{hipopotamo} \textsuperscript{(22)} e \VUgras{dro}\ccplein
\end{VUpage}

% --- Feuillet 104 (folio 102) ----------------------------------------
\begin{VUpage}[104]{102}
%%K t16-03-1 p suite
\VUcontinue
\VUgras{medaro} \textsuperscript{(23)} giboza finis la kolektajo. Me resta\cc
bus ibe dum plu multa hori se ne esabus apude\nl
splendida kampeyo plena de \VUgras{plomba soldati} \textsuperscript{(29)}\nl
qua atraktis mea atenco.

%%K t16-tit-2 sub
\VUsaut{4.0mm}
\VUpk{10.2pt}{40}{Soldati e milito.}
\VUsaut{2.8mm}

%%K t16-04-1 p
\VUblancAlinea
4. --- Mea avulo, qua esas \VUgras{invalida} \textsuperscript{(24)} e qua\nl
devis livar l'armeo pro \VUgras{vunduro} qua igis lu\nl
obtenar la \VUgras{militistala medalio}, tre prizas na\cc
racar a me la \VUgras{kampanii} en qui il partoprenis.\nl
Il esis felica explikante la diversa movadi di\nl
la mikra armeo miniatura. Grupo de vera sol\cc
dati vizitanta la bazaro : \VUgras{jenio-soldato} \textsuperscript{(25)}, \VUgras{Al}\cc
\VUgras{pala-chasero} \textsuperscript{(26)} kapkovrita per \VUgras{beredo},\nl
\VUgras{mayor-serjento} \textsuperscript{(27)} di infantrio e \VUgras{artilrio-ser}\cc
\VUgras{jento} \textsuperscript{(28)} kun ora \VUgras{galono} sur maniko, haltis\nl
proxim ni por askoltar la expliki.

%%K t16-05-1 p
\VUblancAlinea
5. --- Mea avulo, tote fiera, pro ke il havis\nl
tante brilanta audantaro, improvizis kompleta\nl
batalio-plano.

%%K t16-05-2 p
\VUblancAlinea
La fortifikita urbo, kun sua \VUgras{rempari} e \VUgras{fosati},\nl
esis siejata da l'\VUgras{armeo enemika}, qua, pos lon\cc
gatempa \VUgras{bombardado}, esis pronta por \VUgras{asaltar}.\nl
Ma la \VUgras{siejati} koaktata pro la hungro, probabis\nl
\VUgras{ekasaltar} a la \VUgras{kampeyo} dil \VUgras{siejanti}. Pos re\cc
pulsir l'\VUgras{atako} per obstinanta \VUgras{kombato} cirkum\nl
la \VUgras{tendaro}, la \VUgras{vinkinta} siejanti lansis sua \VUgras{eska}\cc
\VUgras{droni} por persequar la \VUgras{vinkita} atakinti. Ici\nl
\VUgras{retretis} kovrante la \VUgras{kombat-agro} per \VUgras{mortinti}\nl
e \VUgras{vunditi}. \VUgras{Kamilisti} forportis ta lasti aden\nl
l'\VUgras{ambulanco} por flegigar li da la \VUgras{kirurgiisto}.

%%K t16-05-3 p
\VUblancAlinea
Malgre la heroala rezisto di kelka \VUgras{batalioni}\nl
qui formacabis \VUgras{quadrato}, la \VUgras{vinkeso} esis kom\ccplein
\end{VUpage}

% --- Feuillet 105 (folio 103) ----------------------------------------
\begin{VUpage}[105]{103}
%%K t16-05-3 p suite
\VUcontinue
pleta, la cetero dil armeo \VUgras{fugis}, livante multa\nl
\VUgras{kaptiti}. L'enemiko iris parfacor sua \VUgras{vinko} oku\cc
pante la urbo. Forsan to esos la fino di la\nl
kampanio e, pos \VUgras{armistico}, esos posibla signa\cc
tar \VUgras{paco-kontrato}.

%%K t16-tit-3 sub
\VUsaut{4.4mm}
\VUpk{10.2pt}{40}{Donacaji por novyaro.}
\VUsaut{2.8mm}

%%K t16-06-1 p
\VUblancAlinea
6. --- La yuna soldati, qui ridis askoltante\nl
la pitoreska naraco dil veterano demandis de\nl
lu kelk altra expliki. Pri me, me cirkumiris la\nl
vendeyo por regardar bel \VUgras{arbalesto} \textsuperscript{(32)} ed\nl
\VUgras{arko} \textsuperscript{(33)} kun lua \VUgras{flecho} \textsuperscript{(31)}. Ibe me trovis altra\nl
kolektajo de bestii : du \VUgras{uceli kantera} \textsuperscript{(34)} : auto\cc
mata \VUgras{naktigalo} e \VUgras{silvio} kantanta per tota fauco,\nl
e serio de stranja bestii : \VUgras{vespertilio} \textsuperscript{(35)},\nl
\VUgras{boao} \textsuperscript{(36)}, \VUgras{tortugo} \textsuperscript{(37)}, \VUgras{mar-hundo} \textsuperscript{(38)}, \VUgras{baleno} \textsuperscript{(39)}\nl
e \VUgras{sharko} \textsuperscript{(40)} bodrucha.

%%K t16-07-1 p
\VUblancAlinea
7. --- Me ne longe haltis por kontemplar li,\nl
nam me videskis lo revata : belega \VUgras{marionet}\cc
\VUgras{teatro} \textsuperscript{(41)} kun splendida \VUgras{dekoruri} e ravisanta\nl
\VUgras{kurteno} reda. Sur la \VUgras{ceno}, du aktori semblis\nl
plear \VUgras{rolo} en tre interesanta \VUgras{teatrajo}, nam la\nl
mikra \VUgras{spektanti} kartona instalita en la lojii o\nl
sidanta sur la \VUgras{stulegi} ed en la \VUgras{partero} dop\nl
l'\VUgras{orkestro-chefo}, aplaudis vivace.

%%K t16-07-2 p
\VUblancAlinea
Regretinde, ta teatro kustis tre chere.

%%K t16-08-1 p
\VUblancAlinea
8. --- Cetere, selektar la ludili esis desfacila,\nl
nam en la vetrini esis amasita omnasorta lu\cc
dili : \VUgras{Arlekino} \textsuperscript{(42)}, \VUgras{Pulchinelo} \textsuperscript{(43)}, \VUgras{Pierrot} \textsuperscript{(44)},\nl
\VUgras{otusi} \textsuperscript{(45)} brandisanta la ali, siflera \VUgras{merli} \textsuperscript{(46)},\nl
aboyera \VUgras{pudeli}, \VUgras{fonografilo} \textsuperscript{(47)} e \VUgras{ludili di pa}\cc
\VUgras{cienteso.} Un de ca ludili tre amuzis me : ol
\parplein
\end{VUpage}

% --- Feuillet 106 (folio 104) ----------------------------------------
\begin{VUpage}[106]{104}
%%K t16-noto-1 noto
\VUnotes{143.2mm}{%
\textit{(*) Videz la tabelo N\textsuperscript{o} 6.}%
}

%%K t16-08-1 p suite
\VUcontinue
reprezentis \VUgras{navala kombato} \textsuperscript{(48)} en qua \VUgras{navo}\nl
esis atakata da \VUgras{pirati} proxim lando plantacita\nl
per \VUgras{kokosieri}.

%%K t16-09-1 p
\VUblancAlinea
9. --- Me povis tre facile kontemplar ta omna\nl
marveli, nam nur poka personi asistis en la\nl
magazino, apene kelka flaneri, \VUgras{dragono} \textsuperscript{(49)} ed\nl
\VUgras{inkasisto} \textsuperscript{(50)} qua prizentis \VUgras{trato} che la chefa\nl
\VUgras{kaso} \textsuperscript{(51)}.

%%K t16-10-1 p
\VUblancAlinea
10. --- Me questionis mea avulo pri l'uzo di\nl
la libri pozita sur tabuli dop la \VUgras{kasistino} ; il\nl
respondis a me ke li esas \VUgras{kontado-libri}.

%%K t16-tit-4 sub
\VUsaut{3.6mm}
\VUpk{10.2pt}{40}{Gapado cirkum butiko.}
\VUsaut{2.8mm}

%%K t16-11-1 p
\VUblancAlinea
11. --- Livante la fako di la ludili, ni adiris\nl
la seliferio donar regardo a la \VUgras{seli} \textsuperscript{(53)}, \VUgras{estri}\cc
\VUgras{bi} \textsuperscript{(54)}, \VUgras{bridi} \textsuperscript{(55)}, \VUgras{morsi} \textsuperscript{(56)} e \VUgras{kravachi}. Dum ke\nl
la \VUgras{fakestro} \textsuperscript{(57)} donis kelk informi a mea avulo,\nl
me iris regardar l'estaleyo di la \VUgras{porcelani} e\nl
\VUgras{kristali} \textsuperscript{(58)} ube esas bela \VUgras{saladuyi} e delikata\nl
\VUgras{tekruchi} \textsuperscript{(59)}.

%%K t16-12-1 p
\VUblancAlinea
12. --- Por irar sur l'unesma etajo, ni pasis\nl
dop la vetrino dil \VUgras{korseti} \textsuperscript{(60)}, apud la kalzo\cc
vendeyo, ube esis estalita tre bela \VUgras{kalzoni} \textsuperscript{(63)}.\nl
Proxime, \VUgras{komizino} \textsuperscript{(61)} igis selektar da \VUgras{kom}\cc
\VUgras{prantino} \textsuperscript{(62)} bel silka \VUgras{texuri} \textsuperscript{(64)}.

%%K t16-12-2 p
\VUblancAlinea
Acensante per l'eskalero, me remarkis ke\nl
la magazino esas dekorita per Japoniana \VUgras{lan}\cc
\VUgras{terni} \textsuperscript{(65)}, \VUgras{parasuni} \textsuperscript{(66)} e grandega \VUgras{palmieri} \textsuperscript{(70)}.

%%K t16-13-1 p
\VUblancAlinea
13 . --- Sur l'unesma etajo, ne esis multo vi\cc
denda; nur mobli (*), \VUgras{ferkesti} \textsuperscript{(67)}, \VUgras{komodi} \textsuperscript{(68)}
\parplein
\end{VUpage}

% --- Feuillet 107 (folio 105) ----------------------------------------
\begin{VUpage}[107]{105}
%%K t16-13-1 p suite
\VUcontinue
\VUgras{infanto-veturi} \textsuperscript{(69)}. Ma l'ensembla vido dil ma\cc
gazino esis tre \VUgras{amuzanta.}

%%K t16-14-1 p
\VUblancAlinea
14. --- Ye nia pedi, me vidis la muzik-ins\cc
trumenti : \VUgras{harpi} \textsuperscript{(71)}, \VUgras{gitari} \textsuperscript{(72)}, \VUgras{mandolini} \textsuperscript{(73)},\nl
\VUgras{valvo-korneti} \textsuperscript{(74)}, \VUgras{klarioneti} \textsuperscript{(75)}, violini kun\nl
lia \VUgras{arketo} \textsuperscript{(76)} e mem \VUgras{tamburego} \textsuperscript{(77)}. Kelke plu\nl
fore, an la paper-fako, \VUgras{studento} \textsuperscript{(78)} marchan\cc
dis de la \VUgras{komizo} \textsuperscript{(79)} kayeruyo. Proxim li, me\nl
distingis bel \VUgras{alfabet-libro} \textsuperscript{(80)}.

%%K t16-14-2 p
\VUblancAlinea
Meze dil magasino esis erektita alta \VUgras{Krist}\cc
\VUgras{nasko-arboro} \textsuperscript{(81)} tote provizita per bujii, kozeti\nl
ed omnasorta ludili : \VUgras{ligna kavali} \textsuperscript{(82)}, \VUgras{pupei} \textsuperscript{(83)},\nl
e. c..

%%K t16-14-3 p
\VUblancAlinea
La \VUgras{parfum-vendeyo} \textsuperscript{(84)} kun sua \VUgras{dentifrici} e\nl
diversa \VUgras{parfumi} difuzis tre bon odoro qua exha\cc
lesis ad ni.

%%K t16-tit-5 sub
\VUsaut{3.4mm}
\VUpk{10.2pt}{40}{Ni kompras kelko.}
\VUsaut{2.8mm}

%%K t16-15-1 p
\VUblancAlinea
15. --- Pro ke mea avulo komencis judikar\nl
tro longa la tempo, ni decensis per la grand\nl
eskalero. Il kelke haltis avan \VUgras{astronomio-ta}\cc
\VUgras{belo} \textsuperscript{(85)} reprezentanta, il dicis a me, \VUgras{kometi} \textsuperscript{(a)},\nl
\VUgras{lun-e sun-eklipsi} \textsuperscript{(b)}, vento-rozo kun la \VUgras{quar}\nl
\VUgras{punti kardinala} \textsuperscript{(c)} \VUgras{(Nordo, Sudo, Esto e Wes}\cc
\VUgras{to)}, mapo di la du \VUgras{mi-sferi} \textsuperscript{d)}, la \VUgras{planeti} \textsuperscript{(e)} e la\nl
\VUgras{sun-sistemo} \textsuperscript{(f)}. Me komprenis nulo en to\nl
omna, e me esis multe plu interesata da la ga\cc
lonizita livreo di \VUgras{garsono livrista} \textsuperscript{(86)} qua pasis,\nl
e dal bela \VUgras{ventizili} \textsuperscript{(87)} qui esis apud me,\nl
proxim la \VUgras{monetuyi} \textsuperscript{(88)} e la \VUgras{portfolii} \textsuperscript{(89)}.

%%K t16-16-1 p
\VUblancAlinea
16. --- Me anke admiris sur la estaleyo \VUgras{plu}\cc
\VUgras{mi} \textsuperscript{(90)}, e \VUgras{zono-bukli} \textsuperscript{(91)}, la bela \VUgras{flori artifica}\cc
\VUgras{la} \textsuperscript{(94)} gracioze aranjita en korbi. La \VUgras{violi}, \VUgras{jil}\ccplein
\end{VUpage}

% --- Feuillet 108 (folio 106) ----------------------------------------
\begin{VUpage}[108]{106}
%%K t16-16-1 p suite
\VUcontinue
\VUgras{flori}, \VUgras{hiacinti} \textsuperscript{(95)}, \VUgras{geranii} \textsuperscript{(96)}, \VUgras{lilii} \textsuperscript{(97)} e \VUgras{kapu}\cc
\VUgras{chini} \textsuperscript{(98)} esis tre bone kopiita.

%%K t16-tit-6 sub
\VUsaut{4.4mm}
\VUpk{10.2pt}{40}{Donacajo de mea avulo.}
\VUsaut{2.8mm}

%%K t16-17-1 p
\VUblancAlinea
17. --- Me pregis mea avulo komprar por me\nl
un de la \VUgras{dento-brosili} \textsuperscript{(92)} qui esis en buxo\nl
apud \VUgras{har-pingli} \textsuperscript{(93)}. Il konsentis e me iris ipsa\nl
pagar mea komprajo an la kaso, proxim qua\nl
on preparis importanta \VUgras{expediajo} \textsuperscript{(100)} por\nl
\VUgras{kliento} \textsuperscript{(99)}.

%%K t16-18-1 p
\VUblancAlinea
18. --- Subite lejera trublo eventis inter la\nl
personi haltinta proxim l'artala \VUgras{bronzaji} \textsuperscript{(101)}.\nl
\VUgras{Furtisto} \textsuperscript{(102)} esis jus kaptita, en sua delikto,\nl
da \VUgras{inspektisto} \textsuperscript{(103)}, dum ke il probis spoliar\nl
ulu. On sorgis venigar policisto por \VUgras{arestar} lu,\nl
e jus kande ni ekiris, la deliktinto esis duktata\nl
aden la karcero.

\VUsaut{6.0mm}
\VUfilet{22mm}
\end{VUpage}
